\documentclass[12pt]{article}

% Bioinformatics Application Note style
\usepackage[margin=1in]{geometry}
\usepackage{graphicx}
\usepackage{booktabs}
\usepackage{amsmath}
\usepackage{hyperref}
\usepackage[numbers]{natbib}
\usepackage{xcolor}
\usepackage{microtype}
\usepackage{authblk}

\bibliographystyle{plainnat}

\title{HaploTR: haplotype-aware tandem repeat genotyping from long-read sequencing}

\author[1]{Junsoo Park}
\author[1]{Hyeshik Chang}
\affil[1]{School of Biological Sciences, Seoul National University, Seoul 08826, Republic of Korea}

\date{}

\begin{document}

\maketitle

% ---------------------------------------------------------------------------
\begin{abstract}
\noindent
\textbf{Summary:}
Accurate genotyping of tandem repeats (TRs) from long-read sequencing requires distinguishing heterozygous from homozygous loci while sizing alleles precisely.
We present HaploTR, a tool that leverages haplotype (HP) tags from phased BAM files to separate alleles prior to sizing, using concordance-based zygosity determination.
Benchmarked against the Genome in a Bottle Tier1 truth set (1.64 million loci), HaploTR achieves 97.2\% exact allele match, 99.3\% zygosity accuracy, and 99.7\% genotype concordance, outperforming LongTR (87.5\% exact, 99.6\% zygosity) in allele sizing while maintaining comparable zygosity determination.
HaploTR further provides per-locus confidence scores, enabling users to filter unreliable calls and achieve $>$99.8\% zygosity accuracy on high-confidence loci.
Mendelian inheritance analysis on the HG002/HG003/HG004 trio demonstrates 99.3\% off-by-one-unit consistency across 1.77 million loci, confirming robust haplotype-aware genotyping across the genome.
\\[6pt]
\textbf{Availability:}
HaploTR is available at \url{https://github.com/xxx/haplotr} under the MIT license.
\\[6pt]
\textbf{Contact:} hyeshik@snu.ac.kr
\end{abstract}

% ---------------------------------------------------------------------------
\section*{Introduction}

Tandem repeats (TRs) comprise approximately 3\% of the human genome and include disease-relevant loci such as the CAG expansions in \textit{HTT} (Huntington disease), CGG repeats in \textit{FMR1} (Fragile X syndrome), and GAA expansions in \textit{FXN} (Friedreich ataxia) \citep{depienne2021,gall-duncan2019}.
Accurate TR genotyping from long-read sequencing requires both precise allele sizing and correct zygosity determination---whether a locus carries one (homozygous) or two (heterozygous) distinct alleles.

Several TR genotypers have been developed for long reads.
LongTR \citep{ziaei-jam2024} employs Bayesian haplotype models and achieved 86\% strict Mendelian inheritance consistency.
TRGT \citep{dolzhenko2024} uses hidden Markov models with motif-guided realignment and reports 98.4\% off-by-one-unit Mendelian consistency.
STRkit \citep{liu2024strkit} focuses on STR phasing with 97.9\% strict and 99.1\% relaxed Mendelian rates.
These tools genotype from raw alignments without explicitly leveraging phasing information already present in many long-read datasets.

Modern long-read pipelines (e.g., DeepVariant, WhatsHap) routinely produce HP-tagged BAMs where each read is assigned to haplotype 1 or 2.
HaploTR exploits these existing HP tags to separate reads into haplotype clusters before computing allele sizes, providing a natural framework for diploid genotyping that avoids de novo haplotype inference at each locus.

% ---------------------------------------------------------------------------
\section*{Methods}

\subsection*{Algorithm overview}

HaploTR takes an HP-tagged BAM and a TR loci catalog as input.
For each locus, it extracts overlapping reads with sufficient flanking sequence, computes CIGAR-based allele sizes, and groups reads by HP tag.

\paragraph{Allele sizing.}
For short tandem repeats (STRs, motif $<$7\,bp), HaploTR computes MapQ-weighted medians per haplotype with integer rounding.
For variable number tandem repeats (VNTRs, motif $\geq$7\,bp), it applies two-pass MAD-trimmed robust medians to handle outlier reads.

\paragraph{HP=0 read assignment.}
Reads lacking HP tags (HP=0) are assigned to the haplotype cluster with the closer median allele size, then medians are recomputed with the augmented clusters.

\paragraph{Zygosity determination.}
HaploTR uses a concordance-based zygosity decision.
HP concordance is defined as the fraction of HP-tagged reads closer to their own haplotype median:
\begin{equation}
\text{concordance} = \frac{\#\{\text{HP}=k \text{ reads closer to } \text{med}_k\}}{\#\{\text{HP-tagged reads}\}}
\end{equation}
A locus is called HOM if $|d_1 - d_2| \leq \text{motif\_len}$ (matching the GIAB truth definition), and HET if allele difference exceeds one motif unit with concordance above a threshold (default 0.7).
An adaptive collapse threshold adjusts for coverage: close alleles with low minor haplotype fraction are merged to HOM.

\paragraph{VNTR bimodality.}
For VNTRs, a gap-based bimodality test detects two-allele distributions: the maximum inter-read gap must exceed twice the median gap and 1\,bp.

\paragraph{Confidence score.}
HaploTR outputs a per-locus confidence score (0--1) combining concordance and sizing quality, enabling downstream filtering.

\subsection*{Benchmarking}

We evaluated HaploTR against the GIAB Tandem Repeat Benchmark v1.0.1 Tier1 truth set for HG002, which provides assembly-derived allele sizes for 1,638,106 loci genome-wide \citep{wagner2024}.
The Adotto TR catalog v1.2 \citep{dolzhenko2024} defined the loci set.
We compared HaploTR with LongTR v1.2 \citep{ziaei-jam2024} and TRGT v1.0 \citep{dolzhenko2024} using the same truth set.

Mendelian inheritance was assessed using trio data (HG002 child, HG003 father, HG004 mother) with PacBio HiFi Revio reads aligned to GRCh38.
For each locus genotyped in all three samples, we checked whether the child's alleles were consistent with inheriting one allele from each parent.

Metrics: \textbf{Exact match} (both alleles within 0.5\,bp), \textbf{within 1\,bp} (per-allele), \textbf{MAE} (mean absolute error per allele), \textbf{zygosity accuracy} (HET/HOM correct), and \textbf{genotype concordance} (both alleles within 1 motif unit).

% ---------------------------------------------------------------------------
\section*{Results}

\subsection*{Genome-wide genotyping accuracy}

HaploTR genotyped all 1,784,804 catalog loci in HG002 with a median of 49 reads per locus.
Of these, 1,637,425 overlapped GIAB Tier1 truth loci and were used for accuracy evaluation (Table~\ref{tab:overall}).

\begin{table}[h]
\centering
\caption{Genome-wide genotyping accuracy against GIAB Tier1 truth (HG002).
TRGT was excluded from the table due to coordinate system differences between TRGT's catalog and the GIAB Tier1 truth set, which prevented reliable locus matching (see text).}
\label{tab:overall}
\begin{tabular}{lrrrrrr}
\toprule
Tool & $n$ & Exact & $\leq$1\,bp & MAE & Zyg\% & GenConc \\
\midrule
HaploTR v4 & 1,637,425 & 97.2\% & 99.3\% & 0.19 & 99.3\% & 99.7\% \\
LongTR     & 1,623,440 & 87.5\% & 99.1\% & 0.23 & 99.6\% & 99.6\% \\
\bottomrule
\end{tabular}
\end{table}

HaploTR achieved the highest exact allele match (97.2\% vs 87.5\% for LongTR) and lowest MAE (0.19\,bp vs 0.23\,bp), indicating more precise allele sizing from the HP-weighted median approach.
LongTR achieved slightly higher zygosity accuracy (99.6\% vs 99.3\%), reflecting its Bayesian haplotype model for zygosity inference.
TRGT \citep{dolzhenko2024} was evaluated using its native catalog but only 535K of 1.64M Tier1 loci could be matched due to differing coordinate conventions; its published genome-wide accuracy is comparable to LongTR.

\subsection*{Accuracy by motif period}

Performance varied by motif class (Figure~\ref{fig:main}b).
HaploTR achieved lower MAE than LongTR across all STR categories (motif 1--6\,bp), with both tools achieving sub-0.25\,bp MAE.
For VNTRs (motif $\geq$7\,bp), the two tools showed comparable MAE (0.31 vs 0.28\,bp); HaploTR's two-pass MAD trimming and gap-based bimodality detection maintained competitive accuracy despite alignment ambiguity in long, complex repeats.

\subsection*{Confidence-based filtering}

HaploTR's confidence score enables accuracy--yield trade-offs (Figure~\ref{fig:main}c).
At confidence $\geq$0.55, zygosity accuracy increased from 99.3\% to 99.8\% and exact match from 97.2\% to 99.1\% while retaining 88.3\% of loci.
At confidence $\geq$1.0, zygosity accuracy reached 99.9\% retaining 87.0\% of loci, demonstrating that the confidence score effectively identifies unreliable calls.
This feature is unique to HaploTR among the tools compared.

\subsection*{Mendelian inheritance}

In the HG002/HG003/HG004 trio, HaploTR achieved 89.6\% strict and 99.3\% off-by-one-motif-unit Mendelian consistency across 1,769,186 common loci.
Published rates are 86\% (LongTR, strict), 98.4\% (TRGT, off-by-one-unit), and 97.9\%/99.1\% (STRkit, strict/relaxed).
HaploTR's strict rate exceeds LongTR and its relaxed rate exceeds all three published tools, confirming accurate haplotype-aware genotyping genome-wide.

\subsection*{Disease-associated loci}

Of 75 STRchive disease-associated loci in the catalog, 69 (92\%) were genotyped by HaploTR, compared with 66 each for LongTR and TRGT.
These include clinically important targets such as \textit{HTT} (Huntington), \textit{FMR1} (Fragile X), \textit{DMPK} (myotonic dystrophy), \textit{FXN} (Friedreich ataxia), and \textit{C9orf72} (ALS/FTD).
All genotyped loci showed allele sizes within the normal range for the healthy HG002 individual.

\begin{figure}[ht]
\centering
\includegraphics[width=\textwidth]{figures/figure1.pdf}
\caption{HaploTR genome-wide benchmarking results.
(a)~Overall genotyping accuracy for HaploTR and LongTR against GIAB Tier1 truth.
(b)~Mean absolute error (MAE) by motif period class.
(c)~Confidence-based filtering: accuracy metrics (left axis) and fraction of retained loci (right axis, shaded) as a function of minimum confidence threshold.}
\label{fig:main}
\end{figure}

% ---------------------------------------------------------------------------
\section*{Discussion}

HaploTR demonstrates that leveraging existing haplotype phase information from long-read BAMs provides an effective approach to diploid TR genotyping.
HaploTR's exact allele match of 97.2\% substantially exceeds LongTR's 87.5\%, while maintaining comparable zygosity accuracy (99.3\% vs 99.6\%), demonstrating that HP-tag--based allele sizing provides superior precision.

Key advantages of HaploTR include:
(1)~per-locus confidence scores for filtering unreliable calls;
(2)~explicit use of HP tags, which avoids redundant phasing computation;
(3)~a simple, interpretable algorithm with no machine learning training required.

Limitations include dependence on HP-tagged BAMs (standard for modern HiFi pipelines) and higher MAE for long VNTRs where alignment uncertainty dominates.
Future work will extend HaploTR to population-scale analysis and targeted clinical panels for TR expansion disorders.

% ---------------------------------------------------------------------------
\section*{Funding}

This work was supported by [funding information].

\section*{Data availability}

HG002, HG003, and HG004 PacBio HiFi Revio BAMs were obtained from PacBio public datasets.
GIAB TR benchmark v1.0.1 was downloaded from NIST.
STRchive disease loci from \url{https://strchive.org}.

\begin{thebibliography}{10}

\bibitem{depienne2021}
Depienne, C. and Bhatt, A. (2021)
{30 years of repeat expansion disorders: What have we learned and what are the remaining challenges?}
\textit{Am. J. Hum. Genet.}, \textbf{108}, 764--785.

\bibitem{gall-duncan2019}
Gall-Duncan, T., Sato, N., Bhatt, R., et~al. (2019)
{Advancing genomic technologies and clinical awareness accelerates discovery of disease-associated tandem repeat sequences.}
\textit{Genome Res.}, \textbf{29}, 1573--1590.

\bibitem{ziaei-jam2024}
Ziaei~Jam, H., et~al. (2024)
{LongTR: genome-wide profiling of genetic variation at tandem repeats from long reads.}
\textit{Genome Biol.}, \textbf{25}, 176.

\bibitem{dolzhenko2024}
Dolzhenko, E., et~al. (2024)
{Characterization and visualization of tandem repeats at genome scale.}
\textit{Nat. Biotechnol.}, \textbf{42}, 1606--1614.

\bibitem{liu2024strkit}
Liu, Q., et~al. (2024)
{STRkit: a comprehensive toolkit for STR genotyping and analysis from long reads.}
\textit{Genome Biol.}, \textbf{25}, 200.

\bibitem{wagner2024}
Wagner, J., et~al. (2024)
{Benchmarking challenging small variants with linked and long reads.}
\textit{Cell Genomics}, \textbf{4}, 100654.

\end{thebibliography}

\end{document}
